% Principal: GDZ PPN236018647, native 23 lower–25, printed 16–18.
\SourceLeaf{23b}\ArithmeticHead{16}{L}{ARITHMETIK.\quad NACHLASS.}
\pdfbookmark[2]{Bemerkungen zu den asymptotischen Gesetzen}{asymptotic-remarks}
\begin{center}\small\bfseries BEMERKUNGEN ZU DEN ASYMPTOTISCHEN GESETZEN DER ZAHLENTHEORIE.\end{center}
\begingroup\small\par\indent
Im Vorstehenden sind diejenigen Aufzeichnungen aus dem Nachlaß zusammengestellt, die sich auf »asymptotische Gesetze« und auf sogenannte »mittlere Werte« der Zahlentheorie beziehen. Die unter [III] [1] und [2] wiedergegebenen Aufzeichnungen gehören dem an anderer Stelle dieses Bandes vollständig abgedruckten »\textit{Tagebuch}« an, die unter [III] [3] wiedergegebene Notiz den gleichfalls weiter unten abgedruckten »\textit{Exercitationes mathematicae}« (1796). Sie sollen aber wegen des sachlichen Zusammenhangs hier zusammen mit den übrigen erläutert werden\footnote{Die folgenden Bemerkungen sind mit Benutzung der im IX. Bericht über den Stand der Herausgabe von \textsc{Gauss}' Werken (Nachrichten der K. Gesellschaft der Wissenschaften zu Göttingen, Geschäftliche Mitteilungen 1911, S. 26 ff.) enthaltenen Ausführungen von P. \textsc{Bachmann} und E. \textsc{Landau}, sowie einiger schriftlichen Mitteilungen von P. \textsc{Bachmann} abgefaßt worden. Die zu benutzenden Zeichen \(O\) und \(o\) haben die folgende Bedeutung (vergl. etwa \textsc{Landau}, Handbuch der Lehre von der Verteilung der Primzahlen I, Leipzig und Berlin 1909, S. 31 und 61): Es seien \(f(x)\) und \(g(x)\) zwei für alle reellen Werte von \(x\), von einem gewissen an, erklärte reelle Funktionen, \(g(x)\) positiv; dann bedeutet \(f(x)=O(g(x))\), daß es zwei Zahlen \(\xi\) und \(A\) gibt, so daß für \(x\geqq\xi\)
\[
|f(x)|<Ag(x)
\]
ist, und \(f(x)=o(g(x))\), daß
\[
\lim_{x\to\infty}\frac{f(x)}{g(x)}=0.
\]
}.
\clearpage
\SourceLeaf{24}\ArithmeticHead{17}{R}{BEMERKUNGEN.}
Den Beweis des Satzes [I] [1], wonach die Anzahl der unter \(a\) liegenden Primzahlen asymptotisch durch \(\frac{a}{\log a}\) dargestellt wird, haben J. \textsc{Hadamard} (Bulletin de la Soc. Mathém. \textbf{24}, 1896, S. 199) und Ch. J. \textsc{de la Vallée Poussin} (Annales de la Société scientif. de Bruxelles \textbf{20}, 2\textsuperscript{e} partie, 1896, S. 360--361) erbracht. Vergl. auch den Werke II, S. 444 ff. abgedruckten Brief von \textsc{Gauss} an \textsc{Encke}.

Der Beweis von [I] [2] ist in einer Arbeit von E. \textsc{Landau} (Bulletin de la Soc. Mathém. \textbf{28}, 1900, S. 25) enthalten, wo gezeigt wird, daß die Menge aller Zahlen \(\leqq a\), die aus \(k\) Primfaktoren zusammengesetzt sind, asymptotisch gleich
\[
\frac1{(k-1)!}\frac{a(\log\log a)^{k-1}}{\log a}
\]
ist.

Die in [I] [3] enthaltene Behauptung hat F. \textsc{Mertens} (\textsc{Crelles} Journal \textbf{78}, 1874, S. 52) unter der Form
\[
\sum_{p\text{ Primzahl}\leqq x}\frac1p
=\log\log x+C-\sum_{m=2}^{\infty}\sum_p\frac1{mp^m}+O\left(\frac1{lx}\right)
\]
bewiesen, wo \(C\) die \textsc{Euler-Mascheroni}sche Konstante bedeutet. Für die von \textsc{Gauss} mit \(V\) bezeichnete Konstante ergäbe sich hiernach der Wert 1,26149, der von dem von \textsc{Gauss} vermuteten etwas abweicht.

Die Behauptung [I] [4] hat L. \textsc{Gegenbauer} (Denkschriften der K. Akademie d. Wiss. Wien \textbf{49}, 1885, Abt. 1, S. 47) bewiesen.

Der Beweis von [I] [5] ergibt sich aus der bei [3] genannten Abhandlung von \textsc{Mertens} (a. a. O. S. 53).

Der Satz [I] [6] ist wohl ein nicht völlig zutreffender Ausspruch der Formel
\[
\sum_{p\leqq x}\left[\frac{x}{p}\right]=x.\log\log x+(V-1)x+o(x),
\]
die aus den Sätzen [I] [1] und [3] hervorgeht.

In [II] [1] ist \(l4\) und \(l5\) im Nenner wohl ein Schreibfehler für \(l5\) und \(l7\), aber selbst dann wäre die Formel nicht ganz zutreffend, da wie \textsc{Gram} (K. Danske Vid. Selsk. Skr. Ser. 6, Bd. 7, 1890--1894, S. 1) gezeigt hat, der richtige Wert der gemeinten Anzahl asymptotisch gleich
\[
\frac1{24}.\frac{(\log n)^4}{\log2.\log3.\log5.\log7}
\]
ist.

Die Bedeutung von [II] [2] haben wir nicht ermitteln können.

Die Aussage [III] [1] hat \textsc{Dirichlet} (Abhandl. der K. preuß. Akad. d. Wissenschaften 1849, Werke \textbf{2}, 1897, S. 59) bewiesen. Bezeichnet man nämlich mit \textsc{Dirichlet} die Summe der Faktoren von \(n\) mit \(f(n)\), so hat nach \textsc{Gauss}
\par\vfill\noindent X\textsubscript{1}.\hfill 3
\clearpage
\SourceLeaf{25}\ArithmeticHead{18}{L}{ARITHMETIK.\quad NACHLASS.}
\[
\frac{f(1)+f(2)+.\,.\,.+f(n)}{\frac12n(n+1)}
\]
\noindent den asymptotischen Wert \(\frac{\pi^2}{6}\), was offenbar auf die von \textsc{Dirichlet} a. a. O. bewiesene Gleichung
\[
f(1)+f(2)+.\,.\,.+f(n)=\frac{\pi^2}{12}n^2+O(n.\log n)
\]
hinauskommt.

Die Aussagen [III] [2] und [3] finden in der von \textsc{Dirichlet} (a. a. O., S. 60--64) bewiesenen Formel
\[
\varphi(1)+\varphi(2)+.\,.\,.+\varphi(n)=\frac3{\pi^2}n^2+O(n^\delta)
\]
ihren Ausdruck, wo \(\varphi(n)\) die Anzahl der zu \(n\) teilerfremden Zahlen nicht größer als \(n\) und \(\delta\) eine zwischen 1 und 2 gelegene Größe bedeutet. In [III] [3] wird der Fall \(n=30\) betrachtet. Es ist für \(n=30\)
\[
\sum_{m=1}^{30}\varphi(m)=278.
\]
\noindent Nun ist aber
\[
(4-1)(3^2-1)(5^2-1)=576=2.288,
\]
\noindent also \textit{angenähert}
\[
\begin{aligned}
\frac{\displaystyle\sum_{m=1}^{30}\varphi(m)}{30.30}
&=\frac12\frac{(2^2-1)(3^2-1)(5^2-1)}{2^2.3^2.5^2}
=\frac12\left(1-\frac1{2^2}\right)\left(1-\frac1{3^2}\right)\left(1-\frac1{5^2}\right)\\
&=\frac12.\frac{1.3}{2.2}.\frac{2.4}{3.3}.\frac{4.6}{5.5}.
\end{aligned}
\]
\noindent Analog ist dann allgemein \textit{angenähert}
\[
\frac{\displaystyle\sum_{m=1}^n\varphi(m)}{nn}
=\frac12\prod_p\left(1-\frac1{p^2}\right)
=\frac12\prod\frac{p-1.p+1}{pp}.
\]

Aus derselben \textsc{Dirichlet}schen Formel läßt sich auch der in [IV] [1] gegebene asymptotische Ausdruck erschließen.

In [IV] [2] finden wir eine Formel für den Mittelwert der Anzahl der Geschlechter binärer quadratischer Formen der Determinante \(N\), die in den artt. 301--302 der \textit{Disqu. Arithm.} in viel schärferer Fassung auftritt.

Der Beweis von [IV] [4] ergibt sich aus der genannten Abhandlung von \textsc{Dirichlet} (a. a. O. S. 52--57); die von \textsc{Gauss} angegebene Konstante stimmt in allen fünf Dezimalen mit dem richtigen Werte \(2C-1\) überein, wo \(C\) die \textsc{Euler-Mascheroni}sche Konstante bedeutet.

Man vergleiche auch die Artikel 26 und 27 des in der zweiten Abteilung dieses Bandes abgedruckten Aufsatzes von P. \textsc{Bachmann} »Über \textsc{Gauss}' zahlentheoretische Arbeiten«.
\begin{flushright}\textsc{Schlesinger}.\end{flushright}
\ClosingRule
\clearpage
\endgroup
